\documentclass[11pt,a4paper]{article}
\usepackage[margin=27mm]{geometry}
\usepackage{amsmath,amssymb}
\usepackage{fontspec,unicode-math}
\setmainfont{STIXGeneral}
\setmonofont{Menlo}
\setmathfont{STIX Two Math}
\usepackage{amsthm,booktabs,longtable,array,enumitem,hyperref,fancyvrb,fvextra}
\hypersetup{colorlinks=true,linkcolor=blue,urlcolor=blue,citecolor=blue}
\newtheorem{theorem}{Theorem}[section]
\newtheorem{lemma}[theorem]{Lemma}
\newtheorem{proposition}[theorem]{Proposition}
\newtheorem{corollary}[theorem]{Corollary}
\theoremstyle{definition}\newtheorem{definition}[theorem]{Definition}
\theoremstyle{remark}\newtheorem{remark}[theorem]{Remark}
\newcommand{\Aut}{\operatorname{Aut}}
\newcommand{\im}{\operatorname{im}}
\newcommand{\id}{\operatorname{id}}
\newcommand{\Stab}{\operatorname{Stab}}
\newcommand{\Ch}{\operatorname{Ch}}
\newcommand{\normal}{\mathrel{\unlhd}}
\setlength{\emergencystretch}{2em}
\date{11 October 2026}

\title{The finite regular-cover computation}
\author{Human Compression project}
\begin{document}\maketitle
This appendix prints the complete mathematical program used in the finite example of \emph{Chirality kernels and the smallest regular cover}. It is copied without alteration from the independently reviewed computation \texttt{phase41\_chiral\_cover\_check.g}. Its logical inputs are the three bare presentations, and its outputs establish all six group rows, the actual subgroup intersections, the mirror and covering maps, the local chirality subgroups, duality and the distinguishing word. GAP's exact algorithms are trusted; this is not a Lean certificate. The main theorem and infinite-family deductions do not depend on this finite computation.

The two old Python runners are preserved separately as operational provenance. Their manuscript-path assumptions are historical, so the current portable \texttt{check.py} executes the unchanged GAP program directly in a fresh directory and checks its complete expected output. Neither an old runner nor a finite validation model supplies an omitted all-parameter proof.
\section*{Complete required program}
\VerbatimInput[fontsize=\scriptsize,breaklines=true,breakanywhere=true,breaksymbolleft={},breaksymbolright={},breakanywheresymbolpre={},breakanywheresymbolpost={}]{programs/chiral_cover.g}

\section*{Finite diagnostics of the reflection formula}
The next complete program checks the two-copy action in three cyclic groups and in the nonabelian symmetric group of degree three. Its countermodels explain the roles of inversion, the copy swap and generation. It is diagnostic evidence only; the quantified theorem is proved in the common paper and its actual action formula is checked in Lean.
\VerbatimInput[fontsize=\scriptsize,breaklines=true,breakanywhere=true,breaksymbolleft={},breaksymbolright={},breakanywheresymbolpre={},breakanywheresymbolpost={}]{programs/reflection_models.py}
\end{document}

